\chapter{The first chapter}

\section{Lorem ipsum}

Lorem ipsum dolor sit amet, consectetur adipiscing elit. Nam magna ex, dictum id convallis vel, tincidunt non lectus.\footnote{This is an example footnote.} Nulla porttitor turpis in felis malesuada, ut semper risus hendrerit.\marginnote{This is an example margin note.} Donec condimentum nulla odio, id aliquet lectus aliquet a.\marginnote{This is an example margin note shifted vertically.}[2cm]

Integer cursus dolor a felis ornare lobortis. Duis varius dolor felis, in elementum nisl vestibulum vel. In sollicitudin justo nisl, sit amet dignissim libero porta hendrerit.\cite{Tianjun_1} Morbi molestie tincidunt justo a lobortis.\cite{McCulloch_Pitts_1}

Etiam ac lectus ac nisl rutrum egestas. Nam ligula nisi, ultricies aliquet ligula ac, sagittis sodales dolor. Sed volutpat lacus at urna convallis, a feugiat augue tincidunt. Nam ex ipsum, placerat ut sem vel, condimentum cursus eros. Donec leo felis, auctor sit amet arcu ut, vehicula congue dolor. Duis augue ex, varius sit amet elit id, blandit facilisis orci. Aliquam euismod eget mi in ullamcorper. Donec hendrerit imperdiet lobortis.

\begin{figure}[tbp]
\centering
\PlaceholderGraphic{0.72\textwidth}{4.8cm}
\caption{Placeholder figure.}
\end{figure}

Duis eu neque non felis aliquet posuere.\MarginPlaceholderFigure{3.2cm}{Margin placeholder figure.} Quisque mollis eros vitae iaculis rhoncus. Nunc at arcu sapien. Aenean aliquam, metus consectetur iaculis elementum, tortor tellus placerat nisi, a condimentum nulla enim et nisi. Maecenas semper volutpat felis tristique bibendum.

Nunc efficitur, nulla et scelerisque euismod, nunc magna pellentesque ligula, quis laoreet justo lorem eu augue. Vestibulum vestibulum, ipsum vitae ultrices commodo, lorem diam semper nisi, non finibus ante justo sed libero. Phasellus molestie maximus nulla id sodales. Integer pulvinar elit ut justo rhoncus cursus. Nullam pharetra dolor vitae mi vulputate, nec congue dui iaculis. Curabitur vehicula molestie mauris, imperdiet semper metus bibendum a. Aliquam vitae est lorem. Cras mollis vel mauris eget interdum. Cras gravida, nibh nec dapibus bibendum, nibh velit imperdiet ligula, in ultrices magna mauris nec orci. 

\MarginFigure[fig:hhlogo]{media/HHLogo-3.png}{HH logo.}{0cm}

Maecenas faucibus dictum libero. Nam pellentesque eros mi, vitae malesuada ex semper a. Suspendisse vitae ultrices diam, sed feugiat mauris. Donec malesuada commodo turpis sit amet consectetur. Mauris posuere id mauris eu accumsan. Cras eu tincidunt magna. Aliquam erat volutpat. Nulla volutpat elit quis lorem vehicula convallis. Sed eu massa sollicitudin dolor eleifend faucibus id non odio. Aenean blandit elit mauris, placerat posuere dui laoreet eu. Nullam feugiat sollicitudin lectus. Curabitur ac tortor ut turpis molestie aliquet. Sed eget varius arcu, placerat imperdiet lorem. Fusce commodo pretium eros, rhoncus egestas velit porttitor non. Phasellus ornare pretium iaculis. See Figure~\ref{fig:hhlogo}.

\newpage

\section[Mathematical symbols $HH\to b\bar{b}\tau^{+}\tau^{-}$]{Mathematical symbols \boldmath ${HH\to b\bar{b}\tau^{+}\tau^{-}}$\unboldmath}

The di-Higgs final state ${HH \to b\bar{b}\tau^{+}\tau^{-}}$ features heavy-flavour jets and leptonic or hadronic ${\tau}$ decays.

Equations without references:

\begin{equation*}
E = m c^{2}\text{.}
\end{equation*}

Equations with references:

\begin{equation}
\int_{0}^{\infty} e^{-x^{2}}\,\mathrm{d}x = \frac{\sqrt{\pi}}{2}\text{.}
\end{equation}

\section{Examples of unit typesetting}

\begin{itemize}
\item Speed of light\index{light!speed of}: ${c=\SI{2.99792458e8}{\metre\per\second}}$
\item Beam energy\index{energy}: ${\qty{6.5}{\TeV}}$ per beam
\item Tracker pitch: ${\qtyrange{50}{75}{\micro\metre}}$
\item Temperature\index{temperature} variation: ${\qty{\pm 0.3}{\kelvin}}$
\item Luminosity: ${\qty{2.0}{\per\femto\barn}}$
\item \SI{12.3}{\kilo\gram}
\item \num{3.45e8}
\end{itemize}

\section{Examples of unit typesetting using electron-volt units with better kerning}
\begin{itemize}
\item Usage in normal text: the beam energy is 6.5~\TeV per beam.
\item Usage in mathematics environment: ${E_{\text{beam}} = 6.5\,\TeV}$.
\item Usage with \texttt{siunitx}: The beam energy is \qty{6.5}{\TeV} per beam.
\item The PETRA accelerator beams at DESY were of \qty{24}{\GeV\per\clight} each.
\end{itemize}

\section{Examples of shorthand commands}

\begin{itemize}
\item Jets are reconstructed using the \antikt clustering algorithm.
\item The \kt algorithm clusters softer constituents first.
\item A Higgs boson can decay to a \bbbar pair.
\item A \bhadron contains at least one \bquark.
\item Secondary vertices can reveal long-lived \bhadrons.
\item The leading \bjet has the highest transverse momentum.
\item Events with two \bjets are retained for further analysis.
\item A \bquark hadronises before it can be observed directly.
\item Two \bquarks are produced in this Higgs-boson decay.
\item The \btagging efficiency is measured in collision data.
\item Each selected jet must satisfy the \btag requirement.
\item The event contains exactly two \btags.
\item A \chadron contains at least one \cquark.
\item Misidentified \chadrons form an important background.
\item \CPT symmetry relates particles to their antiparticles.
\item The \etaphi separation helps distinguish nearby objects.
\item Jets, leptons, photons \etc may be reconstructed in the detector.
\item Charged leptons, \exampligratia electrons, leave tracks in the inner detector.
\item The candidate is neutral, \idest it carries no electric charge.
\item A \fivesigma excess is conventionally associated with a discovery.
\item The \Hbb channel has the largest Higgs-boson branching fraction.
\item The \HT variable is the scalar sum of transverse momenta.
\item The \HThad variable includes only the hadronic contribution.
\item A \largeR jet can contain all the decay products of a boosted particle.
\item \LargeR jets are useful for reconstructing boosted heavy particles.
\item Large \MET can indicate invisible particles in the final state.
\item The selection may be written as ${\MET > \qty{200}{\GeV}}$.
\item The \MTtwo endpoint can constrain the mass scale of a process.
\item The leading-lepton \pT spectrum is compared with the prediction.
\item The \pvalue quantifies the compatibility of the data with a hypothesis.
\item Tracks curve in the \rphi plane when a magnetic field is present.
\item The weak interaction is based on the \SUtwoL gauge group.
\item The \ST requirement suppresses low-energy background events.
\item A \threesigma deviation is evidence, but not normally a discovery.
\item Inclusive \ttbar production is a major background at the LHC.
\item The \ttbarplusbbbar process contains an additional \bquark pair.
\item The \ttbarplusjets sample includes events with additional jets.
\item The \ttbb final state contains two \tquarks and two \bquarks.
\item The \ttplusbb notation highlights \bjets produced with a \tquark pair.
\item The \ttH process directly probes the \tquark Yukawa coupling.
\item The \ttHbb channel combines associated production with a \bquark decay.
\item The \ttHtobb notation makes the Higgs-boson decay explicit.
\item The \tquark is the heaviest elementary particle yet observed.
\item At the LHC, \tquarks are produced both singly and in pairs.
\item A \Wboson mediates charged-current weak interactions.
\item Pairs of \Wbosons provide a sensitive test of electroweak theory.
\item The symbol \Wpm includes both charged states of the weak boson.
\item A \Zboson mediates neutral-current weak interactions.
\item The symbol \Zzero denotes the neutral weak boson.
\end{itemize}
